\documentclass[aspectratio=169,10pt]{beamer}

\usetheme{default}
\usecolortheme{default}
\setbeamertemplate{navigation symbols}{}
\setbeamertemplate{footline}{%
  \hfill\scriptsize Summary Slides \;|\; \insertframenumber/\inserttotalframenumber\hspace{1.0em}\vspace{0.5em}}

\usepackage{amsmath,amssymb,bm}
\usepackage{booktabs}

\definecolor{ngblue}{RGB}{25,64,110}
\definecolor{ngteal}{RGB}{0,114,128}
\definecolor{ngdark}{RGB}{35,35,40}

\setbeamercolor{frametitle}{fg=ngblue}
\setbeamercolor{title}{fg=ngblue}
\setbeamercolor{structure}{fg=ngteal}
\setbeamercolor{normal text}{fg=ngdark,bg=white}
\setbeamerfont{frametitle}{series=\bfseries,size=\large}
\setbeamerfont{title}{series=\bfseries,size=\Large}

\newcommand{\Evec}{\mathbf{E}}
\newcommand{\Dvec}{\mathbf{D}}
\newcommand{\Jvec}{\mathbf{J}}
\newcommand{\rhodef}{\rho_{\mathrm{def}}}
\newcommand{\epssi}{\varepsilon_{\mathrm{Si}}}

\title{Summary Slides}
\subtitle{Baseline, ML, GNN, and superconducting-physics summaries}
\author{}
\date{}

\begin{document}

% -----------------------------------------------------------------------------
\begin{frame}{Module 2 purpose: convert damage into electric field}
\small
\textbf{Module 2 is the electrostatic bridge between radiation-induced defects and carrier transport.}

\vspace{0.45em}
\begin{equation*}
\boxed{
\{C_i,z_i,n,p,N_D^+,N_A^-,\text{bias/BCs}\}
\;\longrightarrow\;
\rho(x,y,t)
\;\longrightarrow\;
\phi(x,y,t)
\;\longrightarrow\;
\Evec(x,y,t)
}
\end{equation*}

\vspace{0.35em}
\begin{itemize}
    \item The central unknown is the electrostatic potential $\phi(x,y,t)$ in a 2-D silicon domain $\Omega$.
    \item Radiation defects enter through their concentrations $C_i(x,y,t)$ and effective charge states $z_i$.
    \item The electric field is not solved independently; it is recovered after the potential solve:
    \[
        \Evec(x,y,t)=-\nabla\phi(x,y,t).
    \]
    \item Module 2 supplies $\phi$ and $\Evec$ to carrier drift-diffusion, recombination, degradation metrics, and later coupled multiphysics updates.
\end{itemize}

\vspace{0.25em}
\textbf{Physical interpretation:} charged defects perturb the local space charge; space charge bends the electrostatic potential; the potential gradient drives carrier motion.
\end{frame}

% -----------------------------------------------------------------------------
\begin{frame}{Governing electrostatics and defect space charge}
\small
\textbf{Electrostatic limit of Maxwell's equations}
\vspace{-0.25em}
\begin{align*}
\nabla\cdot\Dvec &= \rho, & \Dvec &= \epsilon\Evec, & \Evec&=-\nabla\phi,\\[-0.15em]
\Rightarrow\qquad
\nabla\cdot(\epsilon\nabla\phi) &= -\rho.
\end{align*}

\vspace{0.15em}
\textbf{Total semiconductor space charge}
\vspace{-0.15em}
\begin{align*}
\rho(x,y,t)
&= q\big[p(x,y,t)-n(x,y,t)+N_D^+(x,y)-N_A^-(x,y)\big]+\rhodef(x,y,t),\\[-0.05em]
\rhodef(x,y,t)
&= q\sum_{i=1}^{M} z_i C_i(x,y,t).
\end{align*}

\vspace{0.15em}
\begin{itemize}
    \item $p$ and ionized donors $N_D^+$ contribute positive charge; $n$ and ionized acceptors $N_A^-$ contribute negative charge.
    \item A defect species contributes positive, negative, or neutral charge according to the sign and magnitude of $z_i$.
    \item For spatially uniform silicon permittivity,
    \[
        \epsilon=\epssi,
        \qquad
        \epssi \nabla^2\phi=-\rho.
    \]
    \item In nonuniform media or at material interfaces, the conservative form $\nabla\cdot(\epsilon\nabla\phi)=-\rho$ must be retained.
\end{itemize}
\end{frame}

% -----------------------------------------------------------------------------
\begin{frame}{Finite-element weak form and Module 2 outputs}
\footnotesize
\textbf{Weak form on the 2-D domain $\Omega$}
\vspace{-0.25em}
\begin{equation*}
\boxed{
\int_{\Omega} \epsilon\,\nabla w\cdot\nabla\phi\,d\Omega
=
\int_{\Omega} w\rho\,d\Omega
+
\int_{\Gamma_N} w g_N\,d\Gamma
}
\end{equation*}

\vspace{-0.5em}
\begin{align*}
\phi_h(x,y,t) &= \sum_{j=1}^{N} \Phi_j(t)N_j(x,y),\\[-0.15em]
K_{ij} &= \int_{\Omega}\epsilon\,\nabla N_i\cdot\nabla N_j\,d\Omega,\\[-0.15em]
b_i(t) &= \int_{\Omega}N_i\rho(x,y,t)\,d\Omega
+\int_{\Gamma_N}N_i g_N\,d\Gamma,\\[-0.15em]
K\bm{\Phi}(t)&=\bm b(t),
\qquad
\Evec_h|_e(t)=-\nabla\phi_h|_e(t).
\end{align*}

\vspace{-0.15em}
\begin{itemize}
    \item Dirichlet boundaries impose contact voltages or a reference ground: $\phi=\phi_D$ on $\Gamma_D$.
    \item Neumann boundaries impose normal electric displacement or insulating/symmetry conditions: $-\epsilon\nabla\phi\cdot\hat{\mathbf n}=g_N$ on $\Gamma_N$.
    \item Main outputs: nodal potential $\bm\Phi$, reconstructed field $\Evec_h$, and residual diagnostics.
    \item Module 2 is quasi-static: at each coupled time level, the current charge density determines the instantaneous electrostatic solution.
\end{itemize}
\end{frame}



% -----------------------------------------------------------------------------
\begin{frame}{Module 3 purpose: evolve the defect landscape}
\small
\textbf{Module 3 converts the initial radiation damage map into a time-dependent defect concentration field.}

\vspace{0.35em}
\begin{equation*}
\boxed{
\text{Module 1 damage map}
\;\longrightarrow\;
C(x,y,0)\ \text{or}\ S(x,y,t)
\;\longrightarrow\;
C(x,y,t)
\;\longrightarrow\;
\text{Modules 2, 4, and 5}
}
\end{equation*}

\vspace{0.25em}
\textbf{Local defect balance law}
\vspace{-0.25em}
\begin{align*}
\frac{\partial C}{\partial t}+\nabla\cdot\Jvec_C &= R(C,T,\ldots),\\[-0.15em]
\Jvec_C &=-D\nabla C,\\[-0.15em]
\Rightarrow\qquad
\frac{\partial C}{\partial t} &= \nabla\cdot(D\nabla C)+R(C,T,\ldots).
\end{align*}

\vspace{-0.15em}
\begin{itemize}
    \item $C(x,y,t)$ is a coarse-grained concentration of radiation-induced defects in the 2-D silicon domain.
    \item $\Jvec_C$ is the defect flux; Fick's law sends defects down the concentration gradient.
    \item $R$ contains local creation, removal, recombination, clustering, dissociation, or transformation.
    \item For multiple defect species,
    \[
        \frac{\partial C_i}{\partial t}
        =\nabla\cdot(D_i\nabla C_i)+R_i(C_1,C_2,\ldots,T).
    \]
\end{itemize}
\end{frame}

% -----------------------------------------------------------------------------
\begin{frame}{Diffusion, annealing, temperature dependence, and boundaries}
\small
\textbf{Baseline single-species diffusion--reaction model}
\vspace{-0.25em}
\begin{equation*}
\boxed{
\frac{\partial C}{\partial t}
=
\frac{\partial}{\partial x}\!\left(D\frac{\partial C}{\partial x}\right)
+
\frac{\partial}{\partial y}\!\left(D\frac{\partial C}{\partial y}\right)
-k_{\mathrm{ann}}C+S
}
\end{equation*}

\vspace{0.05em}
\textbf{Thermally activated coefficients}
\vspace{-0.25em}
\begin{align*}
D(T)&=D_0\exp\!\left(-\frac{E_D}{k_BT}\right),&
 k_{\mathrm{ann}}(T)&=k_0\exp\!\left(-\frac{E_{\mathrm{ann}}}{k_BT}\right).
\end{align*}

\vspace{-0.15em}
\begin{itemize}
    \item $D$ controls spatial migration; $k_{\mathrm{ann}}$ controls local defect removal or transformation.
    \item Higher $T$ generally increases both migration and annealing by reducing the Boltzmann penalty for crossing activation barriers.
    \item Important timescales:
    \[
        \tau_D\sim\frac{L^2}{D},
        \qquad
        \tau_{\mathrm{ann}}\sim\frac{1}{k_{\mathrm{ann}}}.
    \]
    \item If $\tau_D\ll\tau_{\mathrm{ann}}$, defects redistribute before disappearing; if $\tau_D\gg\tau_{\mathrm{ann}}$, they anneal locally before spreading.
    \item Default closed-domain boundary:
    \[
        \hat{\mathbf n}\cdot\Jvec_C=0,
        \qquad
        \hat{\mathbf n}\cdot D\nabla C=0.
    \]
\end{itemize}
\end{frame}

% -----------------------------------------------------------------------------
\begin{frame}{Finite-element form, update equation, and verification checks}
\footnotesize
\textbf{Weak form for the 2-D defect-evolution equation}
\vspace{-0.25em}
\begin{equation*}
\boxed{
\int_{\Omega} w\frac{\partial C}{\partial t}\,d\Omega
+
\int_{\Omega}D\nabla w\cdot\nabla C\,d\Omega
+
\int_{\Omega}wk_{\mathrm{ann}}C\,d\Omega
=
\int_{\Omega}wS\,d\Omega
}
\end{equation*}

\vspace{-0.35em}
\begin{align*}
C_h(x,y,t)&=\sum_{j=1}^{N}C_j(t)N_j(x,y),\\[-0.15em]
M_{ij}&=\int_{\Omega}N_iN_j\,d\Omega,\\[-0.15em]
(K_D)_{ij}&=\int_{\Omega}D\nabla N_i\cdot\nabla N_j\,d\Omega,\\[-0.15em]
(K_R)_{ij}&=\int_{\Omega}k_{\mathrm{ann}}N_iN_j\,d\Omega.
\end{align*}

\vspace{-0.35em}
\textbf{Backward-Euler update}
\vspace{-0.2em}
\begin{equation*}
\boxed{
\left[M+\Delta t(K_D+K_R)\right]\bm C^{m+1}
=
M\bm C^{m}+\Delta t\,\bm s^{m+1}
}
\end{equation*}

\vspace{-0.35em}
\begin{itemize}
    \item Main output: $C(x,y,t)$, which feeds defect charge in Module 2 and defect-dependent coefficients in later modules.
    \item Pure annealing check: $D=0\Rightarrow C(t)=C(0)e^{-k_{\mathrm{ann}}t}$.
    \item Pure diffusion with zero-flux boundaries: $R=0\Rightarrow M(t)=\int_{\Omega}C\,d\Omega$ should be conserved.
    \item Uniform initial state should remain uniform if sources and boundaries preserve uniformity.
\end{itemize}
\end{frame}



% -----------------------------------------------------------------------------
\begin{frame}{Module 4 purpose: thermal response beyond pure Fourier diffusion}
\small
\textbf{Module 4 computes the 2-D temperature or medium-energy field driven by radiation heating and later coupled electrical heating.}

\vspace{0.25em}
\begin{equation*}
\boxed{
\{\dot q_{\mathrm{rad}},Q_J,Q_{\mathrm{rec}},Q_{\mathrm{def}},\text{thermal BCs}\}
\longrightarrow
T(x,y,t)
\longrightarrow
\{D_i,k_{\mathrm{ann},i},\mu_n,\mu_p,R\}
}
\end{equation*}

\vspace{0.2em}
\textbf{Conservation of thermal energy}
\vspace{-0.25em}
\begin{equation*}
\frac{\partial u}{\partial t}+\nabla\cdot \mathbf q = \dot q,
\qquad
u=C T\quad\text{for constant volumetric heat capacity }C.
\end{equation*}

\vspace{-0.1em}
\begin{itemize}
    \item $u(x,y,t)$ is internal energy density; $T(x,y,t)$ is the equivalent local temperature.
    \item $\dot q$ may represent thermalized radiation energy deposition, Joule heating, recombination heat, or defect-reconfiguration heat.
    \item Module 4 feeds temperature-dependent coefficients back to defect evolution, carrier transport, and coupled degradation metrics.
    \item The motivation for ballistic-diffusive physics is that localized, short-time heating can propagate with finite carrier flight time before reducing to ordinary diffusion.
\end{itemize}
\end{frame}

% -----------------------------------------------------------------------------
\begin{frame}{Fourier, Cattaneo, and ballistic-diffusive thermal transport}
\small
\textbf{Fourier conduction is the local diffusive limit.}
\vspace{-0.25em}
\begin{align*}
\mathbf q_F&=-k\nabla T,\\[-0.15em]
C\frac{\partial T}{\partial t}&=\nabla\cdot(k\nabla T)+\dot q.
\end{align*}

\vspace{-0.15em}
\textbf{Cattaneo adds flux relaxation and finite thermal-wave speed.}
\vspace{-0.25em}
\begin{align*}
\tau\frac{\partial \mathbf q}{\partial t}+\mathbf q&=-k\nabla T,\\[-0.15em]
\tau C\frac{\partial^2 T}{\partial t^2}+C\frac{\partial T}{\partial t}
&=\nabla\cdot(k\nabla T)+\dot q+\tau\frac{\partial \dot q}{\partial t}.
\end{align*}

\vspace{-0.15em}
\begin{itemize}
    \item Fourier is accurate when the mean free path is much smaller than the device length.
    \item Cattaneo fixes the infinite-propagation-speed artifact but is still a local closure.
    \item Ballistic-diffusive modeling splits heat transport into an unscattered ballistic channel and a scattered medium channel:
    \[
        u=u_b+u_m,
        \qquad
        \mathbf q=\mathbf q_b+\mathbf q_m.
    \]
    \item The key dimensionless indicator is the Knudsen number $\mathrm{Kn}=\Lambda/L$.
\end{itemize}
\end{frame}

% -----------------------------------------------------------------------------
\begin{frame}{Module 4 target equation and coupling inputs}
\footnotesize
\textbf{Reduced 2-D ballistic-diffusive medium-temperature equation}
\vspace{-0.25em}
\begin{equation*}
\boxed{
\tau C\frac{\partial^2 T_m}{\partial t^2}
+C\frac{\partial T_m}{\partial t}
=\nabla\cdot(k\nabla T_m)
-\nabla\cdot\mathbf q_b
+\left(\dot q_e+\tau\frac{\partial \dot q_e}{\partial t}\right)
}
\end{equation*}

\vspace{-0.25em}
\begin{equation*}
\dot q_e
=\dot q_{\mathrm{rad}}
+Q_J+Q_{\mathrm{rec}}+Q_{\mathrm{def}}+Q_{\mathrm{ext}}.
\end{equation*}

\vspace{-0.25em}
\begin{itemize}
    \item $\mathbf q_b$ is not a claim that the whole device is ballistic; it represents a retained finite-flight-time correction.
    \item In the diffusive limit, $\tau\rightarrow0$ and $\mathbf q_b\rightarrow0$, the equation reduces to the ordinary heat equation.
    \item Temperature modifies Module 3 defect kinetics:
    \[
        D_i(T)=D_{i0}\exp\!\left(-\frac{E_{D,i}}{k_BT}\right),
        \qquad
        k_{\mathrm{ann},i}(T)=k_{i0}\exp\!\left(-\frac{E_{\mathrm{ann},i}}{k_BT}\right).
    \]
    \item Temperature also modifies carrier mobilities, recombination rates, and later the strength of Joule-heating feedback.
\end{itemize}
\end{frame}

% -----------------------------------------------------------------------------
\begin{frame}{Module 4 finite-element update and verification logic}
\footnotesize
\textbf{Simplified Fourier-limit FEM form used as a baseline thermal path}
\vspace{-0.25em}
\begin{equation*}
\int_{\Omega}w C\frac{\partial T}{\partial t}\,d\Omega
+\int_{\Omega}k\nabla w\cdot\nabla T\,d\Omega
=\int_{\Omega}w\dot q\,d\Omega+\int_{\Gamma_N}wq_N\,d\Gamma.
\end{equation*}

\vspace{-0.4em}
\begin{align*}
T_h&=\sum_j\Theta_j N_j, &
M^T_{ij}&=\int_{\Omega} C N_iN_j\,d\Omega,\\[-0.15em]
K^T_{ij}&=\int_{\Omega} k\nabla N_i\cdot\nabla N_j\,d\Omega, &
\left[M^T+\Delta t K^T\right]\bm\Theta^{m+1}
&=M^T\bm\Theta^m+\Delta t\bm f_T^{m+1}.
\end{align*}

\vspace{-0.3em}
\begin{itemize}
    \item Dirichlet boundaries impose fixed temperature; Neumann boundaries impose heat flux; Robin boundaries impose convective exchange.
    \item The full Module 4 target keeps ballistic-memory terms, but the Fourier-limit solve is the necessary first verification baseline.
    \item Verification checks: uniform equilibrium preservation, diffusive-limit recovery, localized-pulse spreading, and sensitivity to finite-flight-time corrections.
    \item Main output: $T(x,y,t)$ on the same 2-D domain used by Modules 2, 3, 5, and 6.
\end{itemize}
\end{frame}

% -----------------------------------------------------------------------------
\begin{frame}{Module 5 purpose: carrier transport in the damaged device}
\small
\textbf{Module 5 evolves electron and hole densities under electric fields, diffusion, generation, recombination, and defect-modified material properties.}

\vspace{0.25em}
\begin{equation*}
\boxed{
\{\phi,\mathbf E,C_i,T,G,R,\mu_n,\mu_p,D_n,D_p\}
\longrightarrow
n(x,y,t),\ p(x,y,t),\ \mathbf J_n,\ \mathbf J_p
}
\end{equation*}

\vspace{-0.1em}
\textbf{Particle balance laws}
\vspace{-0.25em}
\begin{align*}
\frac{\partial n}{\partial t}&=\frac{1}{q}\nabla\cdot\mathbf J_n+G-R,\\[-0.15em]
\frac{\partial p}{\partial t}&=-\frac{1}{q}\nabla\cdot\mathbf J_p+G-R.
\end{align*}

\vspace{-0.15em}
\begin{itemize}
    \item Electrons and holes have opposite charge, so their current-divergence signs differ in the continuity equations.
    \item $G$ creates electron-hole pairs; $R$ removes them through recombination.
    \item The electric field comes from Module 2, while defect and temperature fields from Modules 3--4 modify the transport coefficients.
    \item Module 5 produces carrier fields that feed back into the Poisson charge density in Module 6.
\end{itemize}
\end{frame}

% -----------------------------------------------------------------------------
\begin{frame}{Drift-diffusion currents and 2-D carrier PDEs}
\small
\textbf{Current-density closures}
\vspace{-0.25em}
\begin{align*}
\mathbf J_n&=q\mu_n n\mathbf E+qD_n\nabla n,\\[-0.15em]
\mathbf J_p&=q\mu_p p\mathbf E-qD_p\nabla p.
\end{align*}

\vspace{-0.15em}
\textbf{Drift-diffusion PDE system}
\vspace{-0.25em}
\begin{align*}
\boxed{\frac{\partial n}{\partial t}
=\nabla\cdot\left(\mu_n n\mathbf E+D_n\nabla n\right)+G-R,}\\[-0.05em]
\boxed{\frac{\partial p}{\partial t}
=-\nabla\cdot\left(\mu_p p\mathbf E-D_p\nabla p\right)+G-R.}
\end{align*}

\vspace{-0.15em}
\begin{itemize}
    \item Drift terms scale with mobility, density, and electric field.
    \item Diffusion terms respond to carrier-density gradients.
    \item In nondegenerate thermal equilibrium, the Einstein relation gives
    \[
        D_n=\mu_n\frac{k_BT}{q},
        \qquad
        D_p=\mu_p\frac{k_BT}{q}.
    \]
    \item Steady-state current response is obtained by setting $\partial n/\partial t=\partial p/\partial t=0$.
\end{itemize}
\end{frame}

% -----------------------------------------------------------------------------
\begin{frame}{How defects degrade mobility and enhance recombination}
\footnotesize
\textbf{Mobility degradation through additional momentum scattering}
\vspace{-0.25em}
\begin{align*}
\mu_n&=\frac{q\tau_{m,n}}{m_n^*},
&
\frac{1}{\tau_{m,n}}&=\frac{1}{\tau_{m,n}^{\mathrm{lat}}(T)}+\frac{1}{\tau_{m,n}^{\mathrm{dop}}}+\\[-0.2em]
&&&\hspace{4.7em}+\sum_j \sigma_{m,n}^{(j)}v_{\mathrm{th},n}C_j,\\[-0.05em]
\mu_p&=\frac{q\tau_{m,p}}{m_p^*},
&
\frac{1}{\tau_{m,p}}&=\frac{1}{\tau_{m,p}^{\mathrm{lat}}(T)}+\frac{1}{\tau_{m,p}^{\mathrm{dop}}}+
\sum_j \sigma_{m,p}^{(j)}v_{\mathrm{th},p}C_j.
\end{align*}

\vspace{-0.3em}
\textbf{Trap-assisted Shockley--Read--Hall recombination}
\vspace{-0.25em}
\begin{equation*}
R_{\mathrm{SRH}}
=
\frac{np-n_i^2}
{\tau_p(n+n_1)+\tau_n(p+p_1)},
\qquad
\frac{1}{\tau_n}=\sum_j\sigma_{n,j}v_{\mathrm{th}}C_j,
\qquad
\frac{1}{\tau_p}=\sum_j\sigma_{p,j}v_{\mathrm{th}}C_j.
\end{equation*}

\vspace{-0.25em}
\begin{itemize}
    \item Increasing defect concentration generally increases scattering and lowers mobility.
    \item Defects also act as trap populations, shortening carrier lifetime and increasing nonradiative recombination.
    \item These closures are the physical path from radiation damage to degraded current response.
\end{itemize}
\end{frame}

% -----------------------------------------------------------------------------
\begin{frame}{Module 5 boundaries, FEM form, and outputs}
\footnotesize
\textbf{Representative boundary conditions}
\vspace{-0.25em}
\begin{align*}
\text{Ohmic contacts:}&\qquad n=n_c,\quad p=p_c,\\[-0.15em]
\text{Insulating boundaries:}&\qquad \mathbf J_n\cdot\hat{\mathbf n}=0,\quad \mathbf J_p\cdot\hat{\mathbf n}=0,\\[-0.15em]
\text{Surface recombination:}&\qquad \frac{1}{q}\mathbf J_n\cdot\hat{\mathbf n}=S_n(n-n_{\mathrm{eq}}),\quad
-\frac{1}{q}\mathbf J_p\cdot\hat{\mathbf n}=S_p(p-p_{\mathrm{eq}}).
\end{align*}

\vspace{-0.35em}
\textbf{Electron weak-form structure}
\vspace{-0.25em}
\begin{equation*}
\int_{\Omega}w\frac{\partial n}{\partial t}\,d\Omega
+
\int_{\Omega}\nabla w\cdot\left(\mu_n n\mathbf E+D_n\nabla n\right)d\Omega
=
\int_{\Omega}w(G-R)\,d\Omega+\text{boundary terms}.
\end{equation*}

\vspace{-0.25em}
\begin{itemize}
    \item The hole weak form is analogous but carries the opposite drift-diffusion sign convention.
    \item Module 5 is nonlinear because current terms contain products such as $\mu_n n\mathbf E$ and recombination depends on $n$ and $p$.
    \item Main outputs: $n$, $p$, $\mathbf J_n$, $\mathbf J_p$, recombination map, generation map, and current-voltage or degradation metrics.
\end{itemize}
\end{frame}

% -----------------------------------------------------------------------------
\begin{frame}{Module 6 purpose: self-consistent multiphysics integration}
\small
\textbf{Module 6 couples the fields from Modules 2--5 on the same 2-D device cross section.}

\vspace{0.2em}
\begin{equation*}
\boxed{
\{C_j\}\leftrightarrow \rho\leftrightarrow \phi,\mathbf E
\leftrightarrow n,p,\mathbf J_n,\mathbf J_p
\leftrightarrow T
\leftrightarrow \{C_j\}
}
\end{equation*}

\vspace{-0.1em}
\begin{itemize}
    \item Module 2 supplies electrostatic potential and electric field from total charge density.
    \item Module 3 supplies defect concentration maps and defect-dependent charge, traps, and kinetic states.
    \item Module 4 supplies temperature and heat-transport response.
    \item Module 5 supplies carrier densities, currents, generation, recombination, and current-related degradation outputs.
    \item The goal is not four separate simulations, but a mutually consistent set of physical maps at each time level.
\end{itemize}
\end{frame}

% -----------------------------------------------------------------------------
\begin{frame}{Coupled governing equations in Module 6}
\footnotesize
\textbf{A compact coupled PDE system}
\vspace{-0.25em}
\begin{align*}
\frac{\partial C_j}{\partial t}
&=\nabla\cdot(D_j(T)\nabla C_j)+R_j(\{C\},n,p,T)+S_j,\\[-0.1em]
\nabla\cdot(\epsilon\nabla\phi)
&=-q\left[p-n+N_D^+-N_A^-+\sum_j z_jC_j\right],\\[-0.1em]
C_T\frac{\partial T}{\partial t}
&=\nabla\cdot(k\nabla T)+Q_{\mathrm{rad}}+Q_J+Q_{\mathrm{rec}}+Q_{\mathrm{def}},\\[-0.1em]
\frac{\partial n}{\partial t}
&=\nabla\cdot(\mu_n n\mathbf E+D_n\nabla n)+G-R,\\[-0.1em]
\frac{\partial p}{\partial t}
&=-\nabla\cdot(\mu_p p\mathbf E-D_p\nabla p)+G-R,\\[-0.1em]
\mathbf E&=-\nabla\phi.
\end{align*}

\vspace{-0.3em}
\begin{itemize}
    \item The equations are coupled through coefficients, sources, charge density, electric field, and heat generation.
    \item The thermal equation shown here is the compact diffusive version; Module 4 may replace it with the reduced ballistic-diffusive form.
\end{itemize}
\end{frame}

% -----------------------------------------------------------------------------
\begin{frame}{Main physical feedback loops and staged coupling levels}
\footnotesize
\textbf{Key feedback mechanisms}
\vspace{-0.25em}
\begin{align*}
\{C_j,z_j\}&\rightarrow \rho_{\mathrm{def}}\rightarrow \phi,\mathbf E,\\[-0.15em]
\mathbf E&\rightarrow \mathbf J_n,\mathbf J_p\rightarrow n,p,\\[-0.15em]
C_j&\rightarrow \mu_n,\mu_p,\tau_n,\tau_p\rightarrow \mathbf J_n,\mathbf J_p,R,\\[-0.15em]
T&\rightarrow D_j,k_{\mathrm{ann},j},\mu_n,\mu_p,R,\\[-0.15em]
\mathbf J\cdot\mathbf E&\rightarrow Q_J\rightarrow T.
\end{align*}

\vspace{-0.3em}
\textbf{Recommended physical-coupling ladder}
\vspace{-0.25em}
\begin{enumerate}
    \item Level 0: one-way field transfer among modules.
    \item Level 1: electrostatic-carrier feedback, $(n,p)\leftrightarrow \rho\leftrightarrow\phi,\mathbf E$.
    \item Level 2: defect-electrostatic-carrier coupling, $C_j\rightarrow\rho_{\mathrm{def}},N_t\rightarrow R_{\mathrm{SRH}}$.
    \item Level 3: thermal feedback through $T$-dependent coefficients and current-induced heat.
    \item Level 4: full self-consistent iteration among defects, electrostatics, carriers, and temperature.
\end{enumerate}
\end{frame}

% -----------------------------------------------------------------------------
\begin{frame}{Module 6 consistency checks and device-response outputs}
\footnotesize
\textbf{Charge consistency}
\vspace{-0.25em}
\begin{equation*}
\rho_{\mathrm{check}}
=q\left[p-n+N_D^+-N_A^-+\sum_j z_jC_j\right],
\qquad
\epsilon_\rho=
\frac{\|\rho_{\mathrm{Poisson}}-\rho_{\mathrm{check}}\|_2}
{\|\rho_{\mathrm{check}}\|_2+\rho_{\mathrm{floor}}}.
\end{equation*}

\vspace{-0.35em}
\textbf{Energy-source bookkeeping}
\vspace{-0.25em}
\begin{equation*}
Q_{\mathrm{tot}}=Q_{\mathrm{rad}}+Q_J+Q_{\mathrm{rec}}+Q_{\mathrm{def}}+Q_{\mathrm{ext}},
\qquad
P_a(t)=\int_\Omega Q_a(x,y,t)\,d\Omega.
\end{equation*}

\vspace{-0.35em}
\textbf{Carrier conservation in an insulating domain}
\vspace{-0.25em}
\begin{equation*}
\frac{d}{dt}\int_\Omega n\,d\Omega=\int_\Omega(G-R)\,d\Omega,
\qquad
\frac{d}{dt}\int_\Omega p\,d\Omega=\int_\Omega(G-R)\,d\Omega.
\end{equation*}

\vspace{-0.25em}
\begin{itemize}
    \item These diagnostics catch stale-field coupling, unit mistakes, and sign errors that can look numerically stable but be physically wrong.
    \item Device-level outputs may include current, leakage, recombination power, temperature rise, defect inventory, and degradation scores.
\end{itemize}
\end{frame}

% -----------------------------------------------------------------------------
\begin{frame}{Module 2 PINN: surrogate mapping for Poisson electrostatics}
\small
\textbf{Goal: learn the Module 2 input-output map while penalizing violations of electrostatics.}

\vspace{0.25em}
\begin{equation*}
\boxed{
\mathcal G_{\theta}:
\left(\mathbf x,\boldsymbol\mu\right)
\longmapsto
\phi_{\theta}(\mathbf x;\boldsymbol\mu),
\qquad
\mathbf E_{\theta}=-\nabla_{\mathbf x}\phi_{\theta}
}
\end{equation*}

\vspace{-0.15em}
\begin{itemize}
    \item $\mathbf x=(x,y)$ is the spatial coordinate; $\boldsymbol\mu$ contains problem parameters such as bias, defect amplitude, material constants, geometry, and source descriptors.
    \item The finite-element solver remains the reference model; the PINN is a fast differentiable surrogate for repeated Module 2 solves.
    \item A purely supervised neural network learns from pairs
    \[
        \left(\rho,\text{BCs},\boldsymbol\mu\right)\mapsto \phi_{\mathrm{FEM}}.
    \]
    \item A physics-informed neural network also sees the governing equation and boundary conditions through its loss function.
\end{itemize}

\vspace{0.15em}
\textbf{Central distinction:} the network is not a new electrostatic law; it is an approximate representation constrained to behave like the Poisson model.
\end{frame}

% -----------------------------------------------------------------------------
\begin{frame}{Module 2 PINN: Poisson residual and boundary losses}
\footnotesize
\textbf{Strong-form residual used at interior collocation points}
\vspace{-0.25em}
\begin{equation*}
\boxed{
\mathcal R_P(\mathbf x;\theta)
=
\nabla\cdot\left(\epsilon\nabla \phi_{\theta}(\mathbf x)\right)+\rho(\mathbf x)
}
\end{equation*}

\vspace{-0.25em}
For uniform silicon permittivity,
\vspace{-0.15em}
\begin{equation*}
\mathcal R_P(\mathbf x;\theta)
=
\epsilon_{\mathrm{Si}}\left(
\frac{\partial^2\phi_{\theta}}{\partial x^2}
+
\frac{\partial^2\phi_{\theta}}{\partial y^2}
\right)+\rho(\mathbf x).
\end{equation*}

\vspace{-0.3em}
\textbf{Representative loss terms}
\vspace{-0.25em}
\begin{align*}
\mathcal L_{\mathrm{PDE}}&=\frac{1}{N_r}\sum_{k=1}^{N_r}\left|\mathcal R_P(\mathbf x_k;\theta)\right|^2,\\[-0.15em]
\mathcal L_D&=\frac{1}{N_D}\sum_{k=1}^{N_D}\left|\phi_{\theta}(\mathbf x_k)-\phi_D(\mathbf x_k)\right|^2,\\[-0.15em]
\mathcal L_N&=\frac{1}{N_N}\sum_{k=1}^{N_N}\left|-\,\epsilon\nabla\phi_{\theta}(\mathbf x_k)\cdot\hat{\mathbf n}-g_N(\mathbf x_k)\right|^2.
\end{align*}
\end{frame}

% -----------------------------------------------------------------------------
\begin{frame}{Module 2 PINN: training objective and data construction}
\footnotesize
\textbf{Combined supervised-plus-physics objective}
\vspace{-0.25em}
\begin{equation*}
\boxed{
\mathcal L_{\mathrm{total}}
=\lambda_s\mathcal L_{\mathrm{data}}
+\lambda_r\mathcal L_{\mathrm{PDE}}
+\lambda_D\mathcal L_D
+\lambda_N\mathcal L_N
+\lambda_{\mathrm{reg}}\mathcal L_{\mathrm{reg}}
}
\end{equation*}

\vspace{-0.15em}
\begin{itemize}
    \item Supervised samples come from trusted FEM cases:
    \[
        \mathcal D_s=\{(\mathbf x_k,\boldsymbol\mu_k,\phi_{\mathrm{FEM},k})\}_{k=1}^{N_s}.
    \]
    \item Interior residual points do not need FEM labels; they train the model to satisfy Poisson's equation between labeled samples.
    \item Boundary points enforce contact voltages, ground reference, insulating conditions, or prescribed normal flux.
    \item Inputs should be nondimensionalized, for example
    \[
        \hat x=x/L,
        \qquad
        \hat\phi=\phi/\phi_0,
        \qquad
        \hat\rho=\rho/\rho_0.
    \]
    \item The loss weights must be scaled so that the largest-magnitude term does not dominate training for numerical rather than physical reasons.
\end{itemize}
\end{frame}

% -----------------------------------------------------------------------------
\begin{frame}{Module 2 PINN: verification, strengths, and limits}
\small
\textbf{A useful PINN surrogate must be judged against the reference electrostatic solve, not only by low training loss.}

\vspace{-0.1em}
\begin{itemize}
    \item Compare potential and field errors:
    \[
        e_{\phi}=\frac{\|\phi_{\theta}-\phi_{\mathrm{FEM}}\|_2}
        {\|\phi_{\mathrm{FEM}}\|_2+\phi_{\mathrm{floor}}},
        \qquad
        e_E=\frac{\|\mathbf E_{\theta}-\mathbf E_{\mathrm{FEM}}\|_2}
        {\|\mathbf E_{\mathrm{FEM}}\|_2+E_{\mathrm{floor}}}.
    \]
    \item Check physically simple cases: zero charge, uniform space charge, linear potential, localized defect charge, and boundary-condition consistency.
    \item Strength: once trained, the surrogate may accelerate parameter sweeps and sensitivity scans.
    \item Limitation: a PINN can still satisfy a weakly weighted residual while violating charge conservation, boundary behavior, or sharp field gradients.
    \item Practical conclusion: keep FEM as the truth model; use the PINN only after systematic out-of-distribution testing.
\end{itemize}
\end{frame}

% -----------------------------------------------------------------------------
\begin{frame}{Module 2 CPINN: why a causal training curriculum is needed}
\small
\textbf{Module 2 is elliptic and quasi-static, so causality does not mean electrostatic time evolution.}

\vspace{0.25em}
\begin{equation*}
\nabla\cdot(\epsilon\nabla\phi)=-\rho,
\qquad
\mathbf E=-\nabla\phi.
\end{equation*}

\vspace{-0.1em}
\begin{itemize}
    \item A single Poisson solve has no intrinsic marching direction in time.
    \item Causal training becomes meaningful only when Module 2 is embedded in an ordered sequence such as dose, irradiation time, defect amplitude, boundary loading, or coupled-iteration continuation.
    \item The ordered coordinate is denoted $\tau$ or $\alpha$:
    \[
        \rho(\mathbf x;\alpha_0),\rho(\mathbf x;\alpha_1),\ldots,
        \rho(\mathbf x;\alpha_K),
        \qquad 0\leq \alpha_0<\alpha_1<\cdots<\alpha_K\leq1.
    \]
    \item The CPINN objective is to learn easy or early states before enforcing later, harder, more highly loaded states.
\end{itemize}

\vspace{0.1em}
\textbf{Interpretation:} Module 2 CPINN is a causality-respecting curriculum for a family of quasi-static Poisson problems.
\end{frame}

% -----------------------------------------------------------------------------
\begin{frame}{Module 2 CPINN: continuation coordinate and staged residuals}
\footnotesize
\textbf{A convenient continuation path from easy to full electrostatic loading}
\vspace{-0.25em}
\begin{equation*}
\boxed{
\rho(\mathbf x;\alpha)
=(1-\alpha)\rho_{\mathrm{easy}}(\mathbf x)
+\alpha\rho_{\mathrm{full}}(\mathbf x),
\qquad 0\leq\alpha\leq1
}
\end{equation*}

\vspace{-0.15em}
\textbf{Residual at each continuation slice}
\vspace{-0.25em}
\begin{equation*}
\mathcal R_P(\mathbf x,\alpha;\theta)
=\nabla\cdot\left(\epsilon\nabla\phi_{\theta}(\mathbf x,\alpha)\right)+\rho(\mathbf x;\alpha).
\end{equation*}

\vspace{-0.2em}
\textbf{Slice loss}
\vspace{-0.25em}
\begin{equation*}
\mathcal L_k
=\frac{1}{N_k}\sum_{m=1}^{N_k}
\left|\mathcal R_P(\mathbf x_m,\alpha_k;\theta)\right|^2
+\mathcal L_{D,k}+\mathcal L_{N,k}+\mathcal L_{\mathrm{data},k}.
\end{equation*}

\vspace{-0.2em}
\begin{itemize}
    \item $\alpha=0$ may represent zero defect charge, weak source amplitude, simple boundary loading, or an earlier dose state.
    \item $\alpha=1$ is the full target electrostatic problem.
    \item Staged training prevents the network from fitting late/high-damage states before it can satisfy simpler earlier states.
\end{itemize}
\end{frame}

% -----------------------------------------------------------------------------
\begin{frame}{Module 2 CPINN: causal weighting of ordered losses}
\footnotesize
\textbf{One causal-PINN strategy weights later residuals by earlier residual success.}

\vspace{-0.2em}
\begin{equation*}
\boxed{
\mathcal L_{\mathrm{causal}}
=\sum_{k=0}^{K} w_k\mathcal L_k,
\qquad
w_k=\exp\left[-\eta\sum_{j=0}^{k-1}\mathcal L_j\right]
}
\end{equation*}

\vspace{-0.15em}
\begin{itemize}
    \item If early slices have large residuals, then $w_k$ suppresses the influence of later slices.
    \item Once early slices are learned, later slices become active because their causal weights increase.
    \item This is not a physical modification of Poisson's equation; it is a training schedule imposed on the learning objective.
    \item The full loss may be written as
    \[
        \mathcal L_{\mathrm{total}}
        =\mathcal L_{\mathrm{causal}}
        +\lambda_D\mathcal L_D
        +\lambda_N\mathcal L_N
        +\lambda_s\mathcal L_{\mathrm{data}}
        +\lambda_{\mathrm{reg}}\mathcal L_{\mathrm{reg}}.
    \]
    \item The parameter $\eta$ controls how aggressively the curriculum enforces earlier states first.
\end{itemize}
\end{frame}

% -----------------------------------------------------------------------------
\begin{frame}{Module 2 CPINN: project use and verification criteria}
\small
\textbf{The CPINN should be accepted only if it improves ordered-state learning without degrading electrostatic fidelity.}

\vspace{-0.1em}
\begin{itemize}
    \item Compare against the ordinary PINN at every ordered slice:
    \[
        e_{\phi}(\alpha_k)=
        \frac{\|\phi_{\theta}(\cdot,\alpha_k)-\phi_{\mathrm{FEM}}(\cdot,\alpha_k)\|_2}
        {\|\phi_{\mathrm{FEM}}(\cdot,\alpha_k)\|_2+\phi_{\mathrm{floor}}}.
    \]
    \item Check whether late/high-damage states improve because early/low-damage states were learned first.
    \item Verify that boundary errors do not grow at large $\alpha$.
    \item Verify that the Poisson residual decreases across the whole continuation path, not only at the final state.
    \item Useful ordered coordinates for this project include dose, defect source amplitude, applied bias, and coupled-iteration loading.
\end{itemize}
\end{frame}

% -----------------------------------------------------------------------------
\begin{frame}{Module 11 GNN: target physical dependency graph}
\small
\textbf{Module 11 learns directed relations among physical variables, not mesh adjacency or numerical solver order.}

\vspace{0.25em}
\begin{equation*}
\boxed{
\mathcal G_{\mathrm{phys}}=(\mathcal V_{\mathrm{phys}},\mathcal E_{\mathrm{phys}}),
\qquad
X_i\rightarrow X_j
\Longleftrightarrow
X_i\ \text{directly enters the law for}\ X_j.
}
\end{equation*}

\vspace{-0.1em}
\begin{itemize}
    \item Nodes include fields, coefficients, sources, and state variables from Modules 2--5:
    \[
        C_j,\rho,\phi,\mathbf E,T,n,p,\mu_n,\mu_p,D_n,D_p,\tau_n,\tau_p,R,\mathbf J_n,\mathbf J_p,Q_J,Q_{\mathrm{rad}}.
    \]
    \item Edges are directed physical dependencies: algebraic, constitutive, source, reaction, transport, or delayed dynamical effects.
    \item The learned graph is allowed to contain cycles because real coupled physics contains feedback.
    \item The goal is edge discovery, mechanism typing, strength estimation, lag classification, and uncertainty quantification.
\end{itemize}
\end{frame}

% -----------------------------------------------------------------------------
\begin{frame}{Module 11 GNN: typed edges and physical meaning}
\footnotesize
\textbf{Each candidate edge carries more information than present/absent.}

\vspace{-0.2em}
\begin{equation*}
\boxed{
X_i\rightarrow X_j
\quad\mapsto\quad
\left(A_{ij},z_{ij},s_{ij},\ell_{ij},\sigma_{ij}\right)
}
\end{equation*}

\vspace{-0.25em}
\begin{itemize}
    \item $A_{ij}$: learned edge strength or probability.
    \item $z_{ij}$: mechanism type, such as algebraic, constitutive, source, transport, reaction, derivative, or delayed dynamical edge.
    \item $s_{ij}$: sign or monotonic trend when physically meaningful.
    \item $\ell_{ij}$: lag class, distinguishing instantaneous relations from delayed responses.
    \item $\sigma_{ij}$: uncertainty or confidence in the inferred edge.
\end{itemize}

\vspace{-0.15em}
\textbf{Examples of target relations}
\vspace{-0.2em}
\begin{align*}
C_j&\rightarrow \rho_{\mathrm{def}}, &
\rho&\rightarrow \phi, &
\phi&\rightarrow \mathbf E,\\[-0.15em]
C_j&\rightarrow \mu_n,\mu_p, &
C_j&\rightarrow \tau_n,\tau_p, &
T&\rightarrow D_j,k_{\mathrm{ann},j},\\[-0.15em]
\mathbf E,n,p&\rightarrow \mathbf J_n,\mathbf J_p, &
\mathbf J\cdot\mathbf E&\rightarrow Q_J, &
Q_J&\rightarrow T.
\end{align*}
\end{frame}

% -----------------------------------------------------------------------------
\begin{frame}{Module 11 GNN: message passing with latent adjacency}
\footnotesize
\textbf{A variable-level GNN updates each physical node using messages from candidate parent variables.}

\vspace{-0.2em}
\begin{align*}
\mathbf m_{ij}^{(\ell)}
&=A_{ij}\,M_{z_{ij}}\!
\left(\mathbf h_i^{(\ell)},\mathbf h_j^{(\ell)},\mathbf e_{ij},\mathbf g\right),\\[-0.1em]
\bar{\mathbf m}_{j}^{(\ell)}
&=\sum_{i\neq j}\mathbf m_{ij}^{(\ell)},\\[-0.1em]
\mathbf h_j^{(\ell+1)}
&=U\!\left(\mathbf h_j^{(\ell)},\bar{\mathbf m}_j^{(\ell)},\mathbf g\right).
\end{align*}

\vspace{-0.25em}
\begin{itemize}
    \item $\mathbf h_j$ is the learned representation of physical variable $X_j$.
    \item $\mathbf e_{ij}$ stores candidate-edge metadata such as allowed mechanism class, dimensional compatibility, and intervention labels.
    \item $\mathbf g$ stores global context such as radiation condition, bias, material identity, temperature scale, and time index.
    \item The adjacency variable $A_{ij}$ can be supervised when the synthetic generator knows the true graph, or treated as a latent variable when only trajectories are available.
    \item Sparsity is essential: without a penalty or prior, the model may explain everything using too many weak edges.
\end{itemize}
\end{frame}

% -----------------------------------------------------------------------------
\begin{frame}{Module 11 GNN: synthetic trajectories and interventions}
\footnotesize
\textbf{Synthetic data must vary both parameter values and graph structure.}

\vspace{-0.25em}
\begin{equation*}
\mathcal D^{(r)}=
\left\{\boldsymbol\theta^{(r)},A^{(r)},Z^{(r)},L^{(r)},X^{(r)}(t),I^{(r)}\right\}.
\end{equation*}

\vspace{-0.2em}
\begin{itemize}
    \item $\boldsymbol\theta^{(r)}$ contains material, source, transport, trapping, thermal, and boundary parameters.
    \item $A^{(r)}$ is the ground-truth physical adjacency used to generate the case.
    \item Mechanism switches create graph-family variation:
    \[
        s_{ij}\in\{0,1\},
        \qquad
        \gamma_{ij}\sim \text{log-uniform or bounded random strength}.
    \]
    \item Matched interventions isolate edges by changing one mechanism while keeping the same random seed:
    \[
        z_j=0,
        \qquad
        \alpha_{\mu,C}=0,
        \qquad
        \alpha_{\tau,C}=0,
        \qquad
        \gamma_J=0,
        \qquad
        T(\mathbf x,t)=T_0.
    \]
    \item Training, validation, and test splits should be by complete physical case, not by time samples from the same trajectory.
\end{itemize}
\end{frame}

% -----------------------------------------------------------------------------
\begin{frame}{Module 11 GNN: from learned graph to physical coupling sequence}
\small
\textbf{The learned graph may contain feedback, so the output is a partial physical ordering, not a single serial algorithm.}

\vspace{-0.1em}
\begin{equation*}
\mathcal G_{\mathrm{phys}}
\longrightarrow
\{\mathcal C_1,\mathcal C_2,\ldots,\mathcal C_K\}
\longrightarrow
\mathcal G_{\mathrm{cond}}.
\end{equation*}

\vspace{-0.15em}
\begin{itemize}
    \item Strongly connected components $\mathcal C_k$ collect mutually dependent feedback groups.
    \item Collapsing each feedback group into one supernode gives the condensation graph $\mathcal G_{\mathrm{cond}}$, which is acyclic.
    \item A topological ordering of $\mathcal G_{\mathrm{cond}}$ gives a physically admissible coupling sequence among influence groups:
    \[
        \mathcal C_{\pi(1)}\prec \mathcal C_{\pi(2)}\prec\cdots\prec \mathcal C_{\pi(K)}.
    \]
    \item Instantaneous algebraic constraints, such as $\rho\rightarrow\phi\rightarrow\mathbf E$, should be separated from delayed dynamical effects, such as $T\rightarrow D_j\rightarrow C_j(t+\Delta t)$.
    \item Success means recovering known synthetic edges, rejecting ablated edges, quantifying edge uncertainty, and identifying regime-dependent feedback groups.
\end{itemize}
\end{frame}




% -----------------------------------------------------------------------------
\begin{frame}{Module 2 superconductivity: electrostatics in hybrid SC devices}
\small
\textbf{The superconducting extension of Module 2 is still electrostatics, but the charge response and boundary conditions change.}

\vspace{0.2em}
\begin{equation*}
\boxed{
\Omega=\Omega_{\mathrm{sub}}\cup\Omega_{\mathrm{ox}}\cup\Omega_{\mathrm{sc}}\cup\Omega_{\mathrm{JJ}},
\qquad
\mathbf E=-\nabla\phi
}
\end{equation*}

\vspace{-0.2em}
\begin{equation*}
\nabla\cdot(\varepsilon\nabla\phi)
=-\rho_{\mathrm{tot}},
\qquad
\rho_{\mathrm{tot}}
=\rho_{\mathrm{def}}+\rho_{\mathrm{trap}}+\rho_{\mathrm{qp}}+\rho_{\mathrm{ind}}+\rho_{\mathrm{ext}}.
\end{equation*}

\vspace{-0.2em}
\begin{itemize}
    \item Substrate and oxide regions store radiation-induced defect charge and trapped charge.
    \item Superconducting films screen static electric fields and are often modeled as equipotential islands.
    \item Quasiparticle charge imbalance can enter as an additional nonequilibrium charge density.
    \item The normal Module 2 source term $q(p-n+N_D^+-N_A^-)+\rho_{\mathrm{def}}$ is replaced by a hybrid-material charge model.
\end{itemize}
\end{frame}

% -----------------------------------------------------------------------------
\begin{frame}{Module 2 superconductivity: screening and equipotential islands}
\footnotesize
\textbf{In a metal or superconductor, static charge perturbations are screened over a very short length scale.}

\vspace{-0.15em}
\begin{align*}
\rho_{\mathrm{ind}} &\approx -\frac{\varepsilon_{\mathrm{sc}}}{\lambda_{\mathrm{scr}}^2}
\left(\phi-\phi_{\mathrm{sc}}\right),\\[-0.05em]
\nabla\cdot(\varepsilon_{\mathrm{sc}}\nabla\phi)
&= -\rho_{\mathrm{ext}} + \frac{\varepsilon_{\mathrm{sc}}}{\lambda_{\mathrm{scr}}^2}
\left(\phi-\phi_{\mathrm{sc}}\right).
\end{align*}

\vspace{-0.2em}
\begin{itemize}
    \item $\lambda_{\mathrm{scr}}$ is an electrostatic screening length, conceptually closer to Thomas-Fermi screening than London magnetic penetration.
    \item In the thin-film static limit, the numerically simpler model is often
    \[
        \phi|_{\Omega_{\mathrm{sc},i}}=V_i,
        \qquad
        \mathbf E_{\parallel}\approx 0 \quad \text{inside the island}.
    \]
    \item The electrochemical potential is the quantity that equilibrates:
    \[
        \mu_{\mathrm{ec}}=\mu_{\mathrm{chem}}-e\phi.
    \]
    \item Radiation-induced local charge is therefore felt mainly through surface charge, island potential shifts, trap occupation, and offset charge.
\end{itemize}
\end{frame}

% -----------------------------------------------------------------------------
\begin{frame}{Module 2 superconductivity: interfaces, capacitance, and offset charge}
\footnotesize
\textbf{The most important device-level output is often not the local field inside the metal, but induced island charge.}

\vspace{-0.2em}
\begin{align*}
\phi_{\mathrm{d}}&=\phi_{\mathrm{sc}} &&\text{on ideal electrostatic interfaces},\\[-0.05em]
\hat{\mathbf n}\cdot(\mathbf D_{\mathrm{d}}-\mathbf D_{\mathrm{sc}})&=\sigma_{\mathrm{int}} &&\text{with interface trapped charge}.
\end{align*}

\vspace{-0.25em}
\begin{equation*}
\boxed{
\mathbf Q=\mathbf C\mathbf V+\mathbf Q_{\mathrm{off}},
\qquad
Q_{\mathrm{off},i}=\int_{\Omega_{\mathrm{d}}}\rho_{\mathrm{rad}}(\mathbf x)\psi_i(\mathbf x)\,d\Omega
}
\end{equation*}

\vspace{-0.2em}
\begin{itemize}
    \item $\mathbf C$ is the capacitance matrix for the superconducting island network.
    \item $\psi_i$ is a weighting potential: solve Laplace's equation with island $i$ at unit voltage and other conductors grounded.
    \item Radiation-generated defects and trapped charges shift $Q_{\mathrm{off}}$, which perturbs charge-sensitive superconducting devices.
    \item For transmon-like devices, offset charge is partially suppressed but still relevant through traps, quasiparticles, and dielectric loss pathways.
\end{itemize}
\end{frame}

% -----------------------------------------------------------------------------
\begin{frame}{Module 2 superconductivity: minimal equations to carry forward}
\footnotesize
\textbf{The first superconducting Module 2 implementation should not overclaim microscopic condensate dynamics.}

\vspace{-0.2em}
\begin{align*}
\nabla\cdot(\varepsilon\nabla\phi)&=-\left(\rho_{\mathrm{def}}+\rho_{\mathrm{trap}}+\rho_{\mathrm{qp}}\right)
&&\text{in dielectric/substrate regions},\\[-0.05em]
\phi&=V_i
&&\text{on superconducting island }i,\\[-0.05em]
\mathbf Q&=\mathbf C\mathbf V+\mathbf Q_{\mathrm{off}},
\qquad
\mathbf E=-\nabla\phi.
\end{align*}

\vspace{-0.25em}
\begin{itemize}
    \item Use superconducting pads as electrostatic conductors first; add finite screening only when needed.
    \item Track oxide/interface traps separately from bulk substrate defect charge.
    \item Treat quasiparticle charge imbalance as a nonequilibrium correction, not as ordinary semiconductor carrier density.
    \item Outputs to later modules: $\phi$, $\mathbf E$, surface/interface fields, $Q_{\mathrm{off}}$, capacitance sensitivity, and trap-field coupling.
\end{itemize}
\end{frame}

% -----------------------------------------------------------------------------
\begin{frame}{Module 3 superconductivity: structural defects are not quasiparticles}
\small
\textbf{Module 3 remains the slow material-state module: structural disorder, traps, TLSs, and junction-barrier defects.}

\vspace{0.2em}
\begin{equation*}
\boxed{
\mathbf C_{\mathrm{SC}}
=\{C_i^{\mathrm{sub}},C_i^{\mathrm{ox}},C_i^{\mathrm{film}},N_t,f_t,N_{\mathrm{TLS}},C_i^{\mathrm{JJ}}\}
}
\end{equation*}

\vspace{-0.1em}
\begin{itemize}
    \item Cryogenic operation separates slow structural variables from fast excitations.
    \item Structural defects and traps can persist for long times; quasiparticles and pair-breaking phonons are fast auxiliary variables handled mainly by Modules 4--5.
    \item The same radiation event can create substrate damage, oxide trapped charge, interface TLS changes, superconducting-film disorder, and junction-barrier defects.
    \item The project should avoid merging all of these into one warm-semiconductor defect concentration.
\end{itemize}
\end{frame}

% -----------------------------------------------------------------------------
\begin{frame}{Module 3 superconductivity: material-resolved diffusion-reaction}
\footnotesize
\textbf{Each material region can carry its own defect species, mobility, reaction network, and boundary behavior.}

\vspace{-0.15em}
\begin{equation*}
\boxed{
\frac{\partial C_i^{(m)}}{\partial t}
=\nabla\cdot\left(D_i^{(m)}(T)\nabla C_i^{(m)}\right)
+R_i^{(m)}(\mathbf C^{(m)},T,n,p)
+S_i^{(m)}
}
\end{equation*}

\vspace{-0.2em}
\begin{align*}
D_i^{(m)}(T)&=D_{i0}^{(m)}\exp\!\left(-\frac{E_{m,i}^{(m)}}{k_BT}\right),\\[-0.05em]
k_{\mathrm{ann},i}^{(m)}(T)&=k_{i0}^{(m)}\exp\!\left(-\frac{E_{\mathrm{ann},i}^{(m)}}{k_BT}\right).
\end{align*}

\vspace{-0.2em}
\begin{itemize}
    \item At millikelvin base temperature, thermally activated diffusion can be nearly frozen.
    \item Radiation can still produce transient local hot spots, nonthermal rearrangement, charge-state changes, and trap occupation changes.
    \item Interface conditions may allow reflection, trapping, sink behavior, or defect conversion at material boundaries.
\end{itemize}
\end{frame}

% -----------------------------------------------------------------------------
\begin{frame}{Module 3 superconductivity: traps, TLSs, and junction defects}
\footnotesize
\textbf{Interface and junction defects are often more device-relevant than bulk displaced atoms.}

\vspace{-0.2em}
\begin{align*}
\rho_{\mathrm{trap}}(\mathbf x,t)&=q\sum_a z_a N_{t,a}(\mathbf x)f_a(\mathbf x,t),\\[-0.05em]
\frac{\partial f_a}{\partial t}
&=c_a^{+}(1-f_a)-c_a^{-}f_a,\\[-0.05em]
P_{\mathrm{TLS}}(E)&\rightarrow \{\tan\delta,\;1/f\ \text{charge noise},\;\text{qubit frequency noise}\}.
\end{align*}

\vspace{-0.25em}
\begin{itemize}
    \item $N_t$ is a trap density; $f_t$ is trap occupation, which controls charge state and electric-field noise.
    \item TLSs are modeled as localized electric dipoles in a heterogeneous polarizability landscape; they couple to microwave electric fields and can produce dielectric loss or fluctuating offset charge.
    \item Josephson-junction barrier defects can affect critical current, tunneling asymmetry, subgap conductance, and switching noise.
    \item Module 3 should pass trap/TLS/barrier state variables to Modules 2, 5, 6, 7, and 8.
\end{itemize}
\end{frame}

% -----------------------------------------------------------------------------
\begin{frame}{Module 3 superconductivity: coupling outputs and first implementation}
\footnotesize
\textbf{The first useful implementation is a region-labeled structural-state update, not a full microscopic defect simulation.}

\vspace{-0.2em}
\begin{equation*}
\boxed{
\mathbf C_{\mathrm{out}}
\longrightarrow
\{\rho_{\mathrm{def}},\rho_{\mathrm{trap}},Q_{\mathrm{off}},N_{\mathrm{TLS}},\delta I_c,\Gamma_{\mathrm{trap}},\tan\delta\}
}
\end{equation*}

\vspace{-0.15em}
\begin{itemize}
    \item To Module 2: charged structural defects and trap occupation produce electrostatic source terms and offset charge.
    \item To Module 4: defect formation, relaxation, and trapping can exchange energy with phonon populations.
    \item To Module 5: traps and disorder modify quasiparticle trapping, poisoning, tunneling, and recombination channels.
    \item Verification targets:
    \[
        C_i(t)=C_i(0)e^{-k_i t}\quad(D_i=0),
        \qquad
        \int_{\Omega_m}C_i\,d\Omega=\text{constant}\quad(R_i=S_i=0).
    \]
    \item State variables should be grouped by material: substrate, oxide/interface, superconducting film, and junction barrier.
\end{itemize}
\end{frame}

% -----------------------------------------------------------------------------
\begin{frame}{Module 4 superconductivity: one temperature field is insufficient}
\small
\textbf{At cryogenic temperatures, deposited radiation energy partitions among several weakly equilibrated channels.}

\vspace{0.15em}
\begin{equation*}
\boxed{
E_{\mathrm{dep}}
\rightarrow
\{U_{\mathrm{sub\ phonon}},\;U_{\mathrm{pb\ phonon}},\;U_{\mathrm{qp}},\;U_{\mathrm{trap}},\;U_{\mathrm{escape}}\}
}
\end{equation*}

\vspace{-0.1em}
\begin{itemize}
    \item A single warm-device heat equation cannot distinguish harmless low-energy phonons from pair-breaking phonons.
    \item Superconducting films have exponentially suppressed electronic heat capacity below $T_c$.
    \item Interfaces can bottleneck energy flow through Kapitza resistance.
    \item Module 4 should represent substrate phonons, pair-breaking phonons, quasiparticle energy, and escape to the package or heat sink at a reduced-model level.
\end{itemize}
\end{frame}

% -----------------------------------------------------------------------------
\begin{frame}{Module 4 superconductivity: pair-breaking threshold and energy partition}
\footnotesize
\textbf{The superconducting gap creates a threshold for converting phonons into quasiparticles.}

\vspace{-0.1em}
\begin{equation*}
\boxed{
\hbar\omega\geq 2\Delta
\quad\Rightarrow\quad
\text{pair-breaking phonon can create two quasiparticles}
}
\end{equation*}

\vspace{-0.15em}
\begin{align*}
\dot U_{\mathrm{pb}}&=\eta_{\mathrm{pb}}\dot E_{\mathrm{dep}}-P_{\mathrm{pair}}-P_{\mathrm{esc}}-P_{\mathrm{down}},\\[-0.05em]
\dot U_{\mathrm{ph}}&=(1-\eta_{\mathrm{pb}})\dot E_{\mathrm{dep}}+P_{\mathrm{down}}-P_{\mathrm{bath}}.
\end{align*}

\vspace{-0.2em}
\begin{itemize}
    \item $\eta_{\mathrm{pb}}$ is an effective fraction of deposited energy placed into pair-breaking-capable phonons.
    \item Downconversion moves high-energy phonons below the pair-breaking threshold.
    \item Escape terms represent phonon loss to substrate, package, backside metal, or external bath.
    \item The relevant output is not only $T$; it is also the time-integrated pair-breaking energy exposure.
\end{itemize}
\end{frame}

% -----------------------------------------------------------------------------
\begin{frame}{Module 4 superconductivity: quasiparticle-phonon kinetics}
\footnotesize
\textbf{A reduced Rothwarf-Taylor structure couples quasiparticle density and pair-breaking phonon population.}

\vspace{-0.2em}
\begin{align*}
\frac{dn_{\mathrm{qp}}}{dt}
&=2\Gamma_{\mathrm{pb}}N_{\mathrm{pb}}
-R_{\mathrm{qp}}n_{\mathrm{qp}}^2
-\Gamma_{\mathrm{trap}}n_{\mathrm{qp}}
+G_{\mathrm{inj}},\\[-0.05em]
\frac{dN_{\mathrm{pb}}}{dt}
&=-\Gamma_{\mathrm{pb}}N_{\mathrm{pb}}
+\frac{1}{2}R_{\mathrm{qp}}n_{\mathrm{qp}}^2
-\Gamma_{\mathrm{esc}}N_{\mathrm{pb}}
+S_{\mathrm{pb}}.
\end{align*}

\vspace{-0.25em}
\begin{itemize}
    \item Pair-breaking phonons increase $n_{\mathrm{qp}}$.
    \item Quasiparticle recombination emits high-energy phonons, which can break more pairs if they do not escape.
    \item Traps and normal-metal quasiparticle sinks remove quasiparticles from active superconducting regions.
    \item Module 4 supplies $N_{\mathrm{pb}}$, phonon energy, and thermal boundary losses to Module 5 and Module 6.
\end{itemize}
\end{frame}

% -----------------------------------------------------------------------------
\begin{frame}{Module 4 superconductivity: boundaries, hierarchy, and metrics}
\footnotesize
\textbf{The practical path is a hierarchy from scalar thermal diffusion to reduced spectral phonon/QP energetics.}

\vspace{-0.15em}
\begin{align*}
-k_1\nabla T_1\cdot\hat{\mathbf n}&=G_K(T_1-T_2) &&\text{Kapitza/interface boundary},\\[-0.05em]
C_{\mathrm{ph}}\frac{\partial T_{\mathrm{ph}}}{\partial t}
&=\nabla\cdot(k_{\mathrm{ph}}\nabla T_{\mathrm{ph}})+P_{\mathrm{dep}}-P_{\mathrm{esc}}-P_{\mathrm{pair}}.
\end{align*}

\vspace{-0.25em}
\begin{itemize}
    \item Level 0: ordinary thermal diffusion with cryogenic material constants.
    \item Level 1: add Kapitza boundary resistance and package/bath coupling.
    \item Level 2: track pair-breaking phonon population or energy.
    \item Level 3: couple pair-breaking phonons to quasiparticle generation/recombination.
    \item Useful outputs: peak temperature, thermal relaxation time, escaped power, pair-breaking fluence, and QP-generation source for Module 5.
\end{itemize}
\end{frame}

% -----------------------------------------------------------------------------
\begin{frame}{Module 5 superconductivity: carrier response changes below $T_c$}
\small
\textbf{Ordinary drift-diffusion is no longer the dominant transport model in superconducting films.}

\vspace{0.15em}
\begin{equation*}
\boxed{
\mathbf J_{\mathrm{tot}}
=\mathbf J_s+\mathbf J_{\mathrm{qp}}+\mathbf J_{\mathrm{leak}}+\mathbf J_{\mathrm{disp}}
}
\end{equation*}

\vspace{-0.15em}
\begin{itemize}
    \item Semiconductor electrons and holes may still matter in substrates, dielectrics, or radiation bursts.
    \item In superconducting films, the condensate carries dissipationless supercurrent and quasiparticles carry dissipative nonequilibrium current.
    \item Below $T_c$, the relevant mobile excitations are Cooper-pair condensate phase and quasiparticle population, not ordinary $n,p$ drift-diffusion.
    \item Module 5 therefore becomes a hybrid carrier-response module: substrate charge, trap occupation, quasiparticles, Josephson transport, and device-response metrics.
\end{itemize}
\end{frame}

% -----------------------------------------------------------------------------
\begin{frame}{Module 5 superconductivity: condensate and Josephson transport}
\footnotesize
\textbf{Supercurrent is governed by phase, not by Ohmic mobility.}

\vspace{-0.15em}
\begin{align*}
\mathbf J_s
&=\frac{2e\hbar}{m^*}|\Psi|^2
\left(\nabla\theta-\frac{2e}{\hbar}\mathbf A\right),\\[-0.05em]
I_s&=I_c\sin\delta,\\[-0.05em]
V&=\frac{\Phi_0}{2\pi}\frac{d\delta}{dt},
\qquad
\Phi_0=\frac{h}{2e}.
\end{align*}

\vspace{-0.2em}
\begin{itemize}
    \item $\theta$ is condensate phase; $\delta$ is the gauge-invariant phase difference across a Josephson junction.
    \item Radiation-induced junction-barrier defects can shift $I_c$, subgap leakage, and tunneling asymmetry.
    \item Capacitive island dynamics enter through charge variables rather than ordinary drift current across the barrier.
    \item In the lowest-fidelity model, Josephson elements are compact circuit relations coupled to field-derived offset charge and QP rates.
\end{itemize}
\end{frame}

% -----------------------------------------------------------------------------
\begin{frame}{Module 5 superconductivity: quasiparticle diffusion, trapping, and poisoning}
\footnotesize
\textbf{Quasiparticles are the main dissipative carrier population inside superconducting films.}

\vspace{-0.15em}
\begin{equation*}
\boxed{
\frac{\partial n_{\mathrm{qp}}}{\partial t}
=\nabla\cdot(D_{\mathrm{qp}}\nabla n_{\mathrm{qp}})
+G_{\mathrm{qp}}
-R_{\mathrm{qp}}n_{\mathrm{qp}}^2
-\frac{n_{\mathrm{qp}}}{\tau_{\mathrm{trap}}}
-\frac{n_{\mathrm{qp}}}{\tau_{\mathrm{esc}}}
}
\end{equation*}

\vspace{-0.2em}
\begin{align*}
\mathbf J_{\mathrm{qp}}&\approx -eD_{\mathrm{qp}}\nabla n_{\mathrm{qp}}+\sigma_{\mathrm{qp}}\big(-\nabla\phi\big),\\[-0.05em]
\Gamma_1^{\mathrm{qp}}&\propto x_{\mathrm{qp}},
\qquad
x_{\mathrm{qp}}\equiv \frac{n_{\mathrm{qp}}}{2N_0\Delta}.
\end{align*}

\vspace{-0.25em}
\begin{itemize}
    \item Pair-breaking phonons from Module 4 provide $G_{\mathrm{qp}}$.
    \item Normal-metal traps and gap-engineered regions reduce active quasiparticle density.
    \item QP tunneling across a junction can produce poisoning, relaxation, parity switching, and excess loss.
\end{itemize}
\end{frame}

% -----------------------------------------------------------------------------
\begin{frame}{Module 5 superconductivity: offset charge and device metrics}
\footnotesize
\textbf{Radiation-generated charge and traps couple to superconducting islands through capacitance and weighting fields.}

\vspace{-0.15em}
\begin{align*}
\mathbf Q&=\mathbf C\mathbf V+\mathbf Q_{\mathrm{off}},\\[-0.05em]
 n_g&=\frac{Q_{\mathrm{off}}}{2e},\\[-0.05em]
Q_{\mathrm{off},i}&=\int \rho_{\mathrm{trap}}(\mathbf x,t)\psi_i(\mathbf x)\,d\Omega.
\end{align*}

\vspace{-0.25em}
\begin{itemize}
    \item $n_g$ is the dimensionless offset charge relevant to island-based superconducting circuits.
    \item Defects alter device response through $Q_{\mathrm{off}}$, $I_c$, dielectric loss, QP density, and QP tunneling rates.
    \item Useful Module 5 outputs include:
    \[
        n_{\mathrm{qp}},\quad x_{\mathrm{qp}},\quad \Gamma_1^{\mathrm{qp}},\quad
        \Gamma_{\mathrm{poison}},\quad\delta I_c,\quad \delta n_g.
    \]
    \item These are device-degradation variables, not merely carrier-density maps.
\end{itemize}
\end{frame}

% -----------------------------------------------------------------------------
\begin{frame}{Module 6 superconductivity: expanded coupled state vector}
\small
\textbf{The superconducting coupling layer must separate electrostatics, structural defects, phonons, quasiparticles, and device response.}

\vspace{0.1em}
\begin{equation*}
\boxed{
\mathbf u_{\mathrm{SC}}
=\{\phi,\mathbf E,Q_{\mathrm{off}},\mathbf C_{\mathrm{def}},N_t,f_t,N_{\mathrm{TLS}},T_{\mathrm{ph}},N_{\mathrm{pb}},n_{\mathrm{qp}},\delta,I_c,\Gamma_1\}
}
\end{equation*}

\vspace{-0.15em}
\begin{itemize}
    \item Module 2 supplies electrostatics, island potentials, capacitance response, and offset charge.
    \item Module 3 supplies structural defects, traps, TLSs, barrier disorder, and disorder-driven parameter shifts.
    \item Module 4 supplies substrate phonon energy, pair-breaking phonons, escape, and thermal relaxation.
    \item Module 5 supplies quasiparticles, Josephson response, QP tunneling, poisoning, and device-rate metrics.
\end{itemize}
\end{frame}

% -----------------------------------------------------------------------------
\begin{frame}{Module 6 superconductivity: radiation-event partition and residuals}
\footnotesize
\textbf{A radiation event must be partitioned into charge, structural damage, phonons, and superconducting excitations.}

\vspace{-0.15em}
\begin{equation*}
\boxed{
E_{\mathrm{dep}}
\rightarrow
\{S_{\mathrm{def}},\rho_{\mathrm{ion}},S_{\mathrm{trap}},S_{\mathrm{pb}},G_{\mathrm{qp}},P_{\mathrm{heat}}\}
}
\end{equation*}

\vspace{-0.15em}
\begin{align*}
\mathcal R_2&:\quad \nabla\cdot(\varepsilon\nabla\phi)+\rho_{\mathrm{def}}+\rho_{\mathrm{trap}}+\rho_{\mathrm{qp}}=0,\\[-0.05em]
\mathcal R_3&:\quad \partial_t\mathbf C_{\mathrm{def}}-\mathcal F_{\mathrm{def}}(\mathbf C_{\mathrm{def}},T_{\mathrm{ph}},f_t)=0,\\[-0.05em]
\mathcal R_4&:\quad \partial_t\{T_{\mathrm{ph}},N_{\mathrm{pb}}\}-\mathcal F_{\mathrm{ph}}(E_{\mathrm{dep}},n_{\mathrm{qp}})=0,\\[-0.05em]
\mathcal R_5&:\quad \partial_t n_{\mathrm{qp}}-\mathcal F_{\mathrm{qp}}(N_{\mathrm{pb}},\mathbf C_{\mathrm{def}},\mathrm{JJ})=0.
\end{align*}
\end{frame}

% -----------------------------------------------------------------------------
\begin{frame}{Module 6 superconductivity: feedback loops and time-scale separation}
\footnotesize
\textbf{The superconducting coupling graph contains fast excitation loops and slow material-state loops.}

\vspace{-0.2em}
\begin{align*}
E_{\mathrm{dep}}&\rightarrow N_{\mathrm{pb}}\rightarrow n_{\mathrm{qp}}\rightarrow \Gamma_1,\Gamma_{\mathrm{poison}},\\[-0.05em]
n_{\mathrm{qp}}&\rightarrow \text{recombination phonons}\rightarrow N_{\mathrm{pb}},\\[-0.05em]
C_{\mathrm{def}},f_t&\rightarrow \rho_{\mathrm{trap}},Q_{\mathrm{off}}\rightarrow \delta n_g,\\[-0.05em]
C_{\mathrm{JJ}}&\rightarrow I_c,\Gamma_{\mathrm{qp}}\rightarrow \text{device response}.
\end{align*}

\vspace{-0.2em}
\begin{itemize}
    \item Fast variables: ionization charge, pair-breaking phonons, quasiparticles, junction tunneling rates.
    \item Intermediate variables: trap occupation, QP trapping, phonon escape, temperature relaxation.
    \item Slow variables: structural defects, interface disorder, TLS population, junction-barrier degradation.
    \item Coupling must respect time scales; forcing all variables into one explicit update can create artificial physics.
\end{itemize}
\end{frame}

% -----------------------------------------------------------------------------
\begin{frame}{Module 6 superconductivity: consistency checks and first algorithm}
\footnotesize
\textbf{The first implementation should prioritize conservation and correct field-transfer bookkeeping.}

\vspace{-0.2em}
\begin{align*}
E_{\mathrm{dep}}&\approx U_{\mathrm{ph}}+U_{\mathrm{pb}}+U_{\mathrm{qp}}+U_{\mathrm{trap}}+E_{\mathrm{esc}}+E_{\mathrm{loss}},\\[-0.05em]
Q_{\mathrm{off},i}&=\int \rho_{\mathrm{trap}}\psi_i\,d\Omega,
\qquad
\rho_{\mathrm{tot}}=\rho_{\mathrm{def}}+\rho_{\mathrm{trap}}+\rho_{\mathrm{qp}}.
\end{align*}

\vspace{-0.25em}
\begin{enumerate}
    \item Partition each radiation event into structural damage, trapped charge, pair-breaking phonons, heat, and QP sources.
    \item Update slow structural/trap state only on appropriate time scales.
    \item Solve electrostatic offset charge from the current trap/defect/QP charge state.
    \item Update phonon-QP kinetics and pass rates to device-response metrics.
    \item Verify energy partition, charge consistency, zero-source equilibrium, and one-way-coupling limits.
\end{enumerate}
\end{frame}

\end{document}

